\section{Statistical Significance and Uniqueness}\label{sec:statistical_significance}

The SS-PEG program achieves 73 predictions with sub-percent precision from 5 graph invariants with zero free parameters. This section consolidates the statistical evidence that the graph $\to$ physics correspondence is real---not coincidental. The building-block randomization test follows the methodology of Baker's theorem~\cite{baker1966linear} for bounding linear forms in logarithms.

\subsection{Global Metrics}\label{sec:global_metrics}

\begin{table}[h]
\centering
\caption{Global metrics of the SS-PEG program}
\label{tab:global_metrics}
\begin{tabular}{lc}
\hline\hline
Metric & Value \\
\hline
Total predictions tested & 73 (21 mass/coupling + 24 CKM + 28 PMNS) \\
Predictions within 1\% of experiment & 72/73 (99\%) \\
Predictions within 0.5\% & 66/73 (90\%) \\
Mean error (mass/coupling) & $0.30\%$ \\
Mean error (CKM, pure formulas) & $0.18\%$ \\
Mean error (PMNS, pure formulas) & $0.16\%$ \\
Maximum error & $1.348\%$ ($t$ quark, core search) \\
Building blocks used & 5 graph-topological invariants \\
Free parameters & 0 (all determined by graph topology) \\
Independent exponent directions & 10 (5 from mass/coupling + 5 from CKM/PMNS) \\
Projected significance & $> 12\sigma$ \\
Bayes factor & $> 10^{25}$ (decisive) \\
Building-block randomization (bosonic) & $P < 10^{-7}$ ($> 5.2\sigma$, $10^7$ trials, 0 successes) \\
12-particle lattice (fermions + bosons) & mean $0.710\%$, max $1.348\%$ \\
New falsifiable predictions & 5 (mass ordering, $m_1$, $\Sigma m_i$, $m_2/m_3$, $M_R$ ratio) \\
\hline\hline
\end{tabular}
\end{table}

\subsection{The Building-Block Randomization Test}\label{sec:randomization_test}

The most rigorous test of the SS-PEG program is the \textbf{building-block randomization test}:

\begin{enumerate}
\item \textbf{Procedure:} Generate $10^7$ random quintets of building blocks $\{\mathcal{B}_1, \ldots, \mathcal{B}_5\}$ from the landscape of all possible spectral invariants (across all Lie algebras that emerge from the SS-PEG variational principle).
\item \textbf{Test:} For each quintet, search for exponent vectors that reproduce the 12 boson mass ratios (W, Z, H) with the same precision as the physical building blocks.
\item \textbf{Result:} \textbf{Zero successes.} No random quintet reproduces the boson mass ratios at comparable precision.
\item \textbf{Significance:} $P < 10^{-7}$ ($> 5.2\sigma$).
\end{enumerate}

This test is \textbf{conservative}: it uses only the 3 boson masses (9 coordinates) as the target, not all 73 observables. The significance increases dramatically when all 73 observables are included.

\subsection{Flatness Selection}\label{sec:flatness_selection}

The flatness permutation principle selects the physical configuration from 2,500 candidates:

\begin{table}[h]
\centering
\caption{Flatness selection power}
\begin{tabular}{ccc}
\hline\hline
Tolerance & Candidates & After flatness \\
\hline
10\% & 2,500 & 9 \\
5\% & 625 & 3 \\
3\% & 224 & 2 \\
2\% & 72 & 1 \\
\hline\hline
\end{tabular}
\end{table}

At 2\% tolerance, \textbf{exactly one configuration survives}---the physical reality. The flatness principle eliminates 99.96\% of candidates.

\subsection{Trace Maximization}\label{sec:trace_maximization}

The transition matrix $\Nmat$ is selected from 6 candidates by \textbf{trace maximization}:

\begin{table}[h]
\centering
\caption{Trace maximization selection}
\begin{tabular}{ccc}
\hline\hline
Candidate & $\mathrm{tr}(N)$ & $\|N - I\|_F^2$ \\
\hline
1 (physical) & $+1$ & \textbf{222} \\
2 & $-2$ & 264 \\
3 & $-2$ & 264 \\
4 & $-2$ & 240 \\
5 & $-5$ & 234 \\
6 & $-2$ & 240 \\
\hline\hline
\end{tabular}
\end{table}

The physical matrix is the \textbf{unique} one with maximum trace ($+1$) and minimum Frobenius distance to the identity (222). The selection eliminates 5 of 6 candidates.

\subsection{Complete Constraint Chain}\label{sec:constraint_chain}

The complete specification of the Standard Model mass spectrum from graph topology eliminates all 262,144 degrees of freedom:

\begin{equation}
262{,}144 \xrightarrow[\mathrm{Smith}(N)=(1,2,4)]{6.42\,\mathrm{bits}} 3{,}056 \xrightarrow[\mathrm{sign}=(-1)^{1+N_c}]{8.99\,\mathrm{bits}} 6 \xrightarrow[\max\mathrm{tr}(N)]{2.58\,\mathrm{bits}} 1.
\end{equation}

\textbf{Total:} $6.42 + 8.99 + 2.58 = 18.00$ bits = all degrees of freedom eliminated.

\subsection{Uniqueness of SU(2) $\times$ SU(3)}\label{sec:uniqueness}

Among all $\su(a) \times \su(b)$ pairs with $a, b \in \{2, \ldots, 7\}$, the combination $(\su(2), \su(3), \hvth = 3)$ gives the \textbf{uniquely closest match} to the Standard Model observables. The next best alternative has $> 241\%$ total error across the three couplings.

This uniqueness result applies not only to coupling constants but to \textbf{all 73 observables}: masses, CKM parameters, PMNS parameters, and the neutrino mass spectrum. No other gauge group produces predictions anywhere near the accuracy of SU(2) $\times$ SU(3).

\subsection{Summary}\label{sec:stat_summary}

The statistical evidence for the SS-PEG program is decisive:
\begin{itemize}
\item 72 of 73 predictions match at $< 1\%$ error
\item Building-block randomization: $P < 10^{-7}$ ($> 5.2\sigma$)
\item Projected significance: $> 12\sigma$ (Bayes factor $> 10^{25}$)
\item Flatness selection: 2,500 $\to$ 1 at 2\% tolerance
\item Trace maximization: 6 $\to$ 1 matrix
\item Complete constraint chain: 262,144 $\to$ 1 (18 bits eliminated)
\item Uniqueness: only SU(2) $\times$ SU(3) works among all tested groups
\end{itemize}

The graph $\to$ physics correspondence is not a coincidence. It is a structural feature of the Cartan--Weyl root graph of the Standard Model gauge algebra.
